\documentclass[11pt]{article}
\usepackage[margin=1in]{geometry}
\usepackage{amsmath,amssymb,amsthm,booktabs,hyperref}
\hypersetup{colorlinks=true,linkcolor=blue,urlcolor=blue}
\newtheorem{theorem}{Theorem}
\newcommand{\Tr}{\operatorname{Tr}}
\newcommand{\End}{\operatorname{End}}
\newcommand{\E}{\mathbb E}
\newcommand{\ip}[2]{\langle #1,#2\rangle}
\title{An exact obstruction to pairwise orthogonality\\MathDB 376318}
\author{Research note}
\date{17 September 2026}
\begin{document}
\maketitle

\section*{Task 1. Resolve the original conjecture}
Kuijper's Conjecture 2.25 \cite[p.~26]{K} asserts that the multivariate
Hermitian Jacobi polynomials in Definition 2.23 are orthogonal whenever
their indices are non-equivalent. That definition fixes a multivariate
Schur leading term and subtracts strictly lower terms in the specified
partial order; it does not impose orthogonality between incomparable
indices. We disprove the conjecture, including its allowance for
arbitrary choices within repeated representation types.

Set $m=3$, $k=6$, and
\[
 \alpha_j=\beta_j=\tfrac32\quad(1\le j\le6),\qquad
 \tau_1=\cdots=\tau_6=(1).
\]
Thus $X_1,\ldots,X_6$ are independent and uniformly distributed on
$0<X<I_3$ in the real vector space of complex Hermitian matrices.
These parameters lie strictly inside the source's admissible range.
All inner products below are integration against this probability law.
The two final partition labels are
\[
 \lambda=(4,1,1),\qquad \mu=(3,3).
\]
They are distinct and incomparable in dominance order.

\begin{theorem}
For these parameters, the spaces spanned by the source's Jacobi
polynomials with final labels $\lambda$ and $\mu$ are not orthogonal.
This holds for every choice of the multivariate Schur bases within
their multiplicity spaces. Explicit real polynomials $p_\lambda$ and
$p_\mu$ in those respective spans satisfy
\[
 \boxed{\ip{p_\lambda}{p_\mu}
 =-\frac{395}{12107408}\left(\frac3{70}\right)^6
 =-\frac{57591}{284884888758400000}<0.}
\]
Consequently at least one pair of the original Jacobi basis
polynomials with these non-equivalent partition labels has nonzero
inner product, contradicting Conjecture 2.25.
\end{theorem}

\subsection*{1. Matrix moments}
Write $Z=X-I_3/2$. Unitary invariance and invariance under $X\mapsto I-X$
give $\E X=I_3/2$ and
\[
 \E(X_{ab}X_{cd})=A\delta_{ab}\delta_{cd}
                         +B\delta_{ad}\delta_{bc}.
\]
The unordered eigenvalues have density proportional to
$\Delta(x)^2=\prod_{i<j}(x_i-x_j)^2$ on $[0,1]^3$.
Direct integration of this polynomial gives
\[
\begin{array}{c|cccc}
 q(x)&1&\sum_i x_i&\sum_i x_i^2&(\sum_i x_i)^2\\ \hline
 \int_{[0,1]^3}q(x)\Delta(x)^2\,dx
 &1/360&1/240&13/4200&9/1400.
\end{array}
\]
These identities can be checked simply by integrating each monomial,
$\int_{[0,1]^3}x_1^ix_2^jx_3^\ell\,dx=
((i+1)(j+1)(\ell+1))^{-1}$.
Hence $3A+9B=39/35$ and $9A+3B=81/35$, so that
$A=17/70$, $B=3/70$. In particular,
\begin{equation}\label{cov}
 \E(Z_{ab}Z_{cd})
 =a\delta_{ad}\delta_{bc}+b\delta_{ab}\delta_{cd},
 \qquad a=\frac3{70},\quad b=-\frac1{140},\quad b/a=-\frac16.
\end{equation}
This also follows from stationarity of the source's equation (8):
its degree-two moment equations are $12A+2B=3$, $2A+12B=1$.

\subsection*{2. Concrete polynomials and their partition blocks}
Let $V=\mathbb C^3$. Denote by $T_\sigma$ the permutation of tensor
factors on $V^{\otimes6}$ corresponding to $\sigma\in S_6$.
For a partition $\eta\vdash6$ of length at most three, let $f^\eta$
be the dimension of its symmetric-group representation and let
$D_\eta=\dim S_\eta(V)$. The central orthogonal projector onto
the $\eta$-isotypic summand is
\begin{equation}\label{projector}
 E_\eta=\frac{f^\eta}{720}
              \sum_{\sigma\in S_6}\chi^\eta(\sigma)T_\sigma.
\end{equation}
Define the explicit multilinear polynomials
\begin{align*}
 q_\eta(X_1,\ldots,X_6)
   &=D_\eta^{-1}\Tr\bigl(E_\eta(X_1\otimes\cdots\otimes X_6)\bigr),\\
 f_\eta(X_1,\ldots,X_6)
   &=D_\eta^{-1}\Tr\bigl(E_\eta(Z_1\otimes\cdots\otimes Z_6)\bigr).
\end{align*}
They are real on Hermitian arguments and invariant under simultaneous
unitary conjugation and under all permutations of the six arguments.
Formula \eqref{projector} makes them finite, fully specified polynomials.

Here is the identification with the source's multiplicity spaces.
Schur--Weyl decomposition gives
\[
 V^{\otimes6}=\bigoplus_{\eta\vdash6,\,\ell(\eta)\le3}
              S_\eta(V)\otimes F_\eta,
 \qquad
 \End_{U(3)}(V^{\otimes6})=
              \bigoplus_\eta I_{S_\eta(V)}\otimes\End(F_\eta).
\]
The trace pairing identifies this commutant with simultaneous-unitary
invariant multilinear polynomials. Its $\eta$ block, denoted $W_\eta$,
is exactly the direct sum over equivalent labels in source equation
(12), with all six individual partitions equal to $(1)$.
In particular, $q_\eta\in W_\eta$.
We use the whole block: we do not assume that $q_\eta$ is one of the
source's individually normalized Schur basis vectors.

\subsection*{3. All strictly lower terms are accounted for}
Let $L$ be the span of source Schur polynomials whose tuple of
individual partitions is strictly below $((1),\ldots,(1))$ in the
source's order on pages 19--20. At the last differing position $j$,
the partition must be zero. Thus each such polynomial is independent
of $X_j$, even if it has arbitrarily high degree in earlier variables.
Every polynomial multilinear in $Z_1,\ldots,Z_6$ has conditional mean
zero in each variable, and is therefore orthogonal to \emph{all} of $L$.

Let $Cq=q(X_1-I/2,\ldots,X_6-I/2)$ for a multilinear $q$.
Then $Cq-q\in L$ and $Cq\perp L$; in particular $Cq_\eta=f_\eta$.
The only length-at-most-three partitions of six strictly below either
$\lambda$ or $\mu$ are
\[
 \nu=(3,2,1),\qquad \omega=(2,2,2).
\]
Indeed, all such partitions are
$(6),(5,1),(4,2),(4,1,1),(3,3),(3,2,1),(2,2,2)$, and comparison of
their partial sums proves the assertion. Consequently the full
strictly lower space for either target block is the orthogonal sum
\begin{equation}\label{lower}
 L\mathbin{\oplus}N,\qquad N=C(W_\nu\oplus W_\omega).
\end{equation}
This identity uses $w-Cw\in L$ for each $w$ in the two lower blocks.
It does not require a finite-dimensional assumption about $L$.

The product measure is invariant under permutation of the six slots.
The finite-dimensional space $N$ has the same invariance.
In a block $W_\eta$, the slot action is conjugation by the irreducible
$S_6$ representation on $F_\eta$. By Schur's lemma its fixed space
consists precisely of scalar multiples of the identity of $F_\eta$.
Thus the permutation-fixed subspace of $N$ is
\[
 N^{S_6}=\operatorname{span}\{f_\nu,f_\omega\}.
\]
Orthogonal projection onto $N$ commutes with slot permutations.
Since $f_\lambda$ and $f_\mu$ are fixed, their projections belong to
this two-dimensional subspace. Projecting away these two vectors
therefore removes \emph{every} strictly lower term required by the
source, including all multiplicity copies.

For precision, the map taking a leading polynomial $q\in W_\eta$
to the polynomial with that leading term, lower corrections, and
orthogonality to the lower space is linear and unique. Existence here
follows by first applying $C$ and then projecting away $N$.
For uniqueness, the difference of two possible answers is lower and
orthogonal to itself, hence zero because the measure has positive
density on an open set. Applying this linear map to a Schur basis of
$W_\eta$ gives exactly Definition 2.23. Thus its value on $q_\eta$
is a linear combination of the source's Jacobi polynomials with
final label $\eta$, for \emph{any} allowed choice of basis.

\subsection*{4. Exact Gram calculation}
For $r\ge0$, let $f^{\eta/\xi}$ denote the number of standard tableaux
of skew shape $\eta/\xi$, taken as zero unless $\xi\subseteq\eta$.
Write $f^\xi=f^{\xi/\varnothing}$ and $D_\varnothing=1$.
The partial trace of $E_\eta$ over $6-r$ slots is
\begin{equation}\label{partial}
 \Tr_{6-r}E_\eta
  =\sum_{\xi\vdash r,\,\ell(\xi)\le3}
       \frac{D_\eta}{D_\xi}f^{\eta/\xi}E_\xi^{(r)}.
\end{equation}
To see this, it commutes with both $U(3)$ and $S_r$, so is scalar
on each $\xi$ summand. The branching rule from $S_6$ to $S_r$ gives
\[
 \Tr\bigl(E_\eta(E_\xi^{(r)}\otimes I)\bigr)
       =D_\eta f^\xi f^{\eta/\xi}.
\]
Dividing by $\Tr E_\xi^{(r)}=D_\xi f^\xi$ proves \eqref{partial}.

Independence and \eqref{cov} give the following normalized Gram matrix:
\begin{equation}\label{gram}
 G_{\eta\zeta}:=a^{-6}\ip{f_\eta}{f_\zeta}
 =\sum_{r=0}^6\binom6r\left(-\frac16\right)^{6-r}
    \sum_{\xi\vdash r,\,\ell(\xi)\le3}
        \frac{f^\xi}{D_\xi}f^{\eta/\xi}f^{\zeta/\xi}.
\end{equation}
Indeed, choosing the $a$ term of the covariance in $r$ slots and the
$b$ term in the others contracts the two partial traces in
\eqref{partial}. The trace of their product, divided by
$D_\eta D_\zeta$, is the inner sum in \eqref{gram}.
Slot symmetry gives the binomial coefficient.

In the order $(\lambda,\mu,\nu,\omega)$, evaluation gives
\begin{equation}\label{matrix}
 G=\frac1{46656}
 \begin{pmatrix}
 15604&2510&-14576&11330\\
 2510&14647&-4480&3505\\
 -14576&-4480&29152&-42640\\
 11330&3505&-42640&174505
 \end{pmatrix}.
\end{equation}
All entries can be calculated using only the corner-removal recurrence
\[
 f^{\eta/\xi}=
 \sum_{\substack{\eta^-\text{ obtained by removing a corner}\\
                         \xi\subseteq\eta^-}} f^{\eta^-/\xi},
 \qquad f^{\xi/\xi}=1,
\]
and the dimension formula
\[
 D_\xi=\prod_{1\le i<j\le3}\frac{\xi_i-\xi_j+j-i}{j-i},
\]
with zero parts appended. The accompanying exact verifier implements
these formulas. It also verifies all sixteen entries independently,
by constructing the four $729\times729$ integer matrices
$\sum_{\sigma\in S_6}\chi^\eta(\sigma)T_\sigma$ directly and taking
their partial traces. In that calculation characters are computed as
the coefficient of $x^{\eta+(2,1,0)}$ in
\[
 \prod_{i<j}(x_i-x_j)\prod_{h\in\mathrm{cycles}(\sigma)}
          (x_1^{|h|}+x_2^{|h|}+x_3^{|h|}),
\]
without using the skew-tableau recurrence. Every operation is exact.

\subsection*{5. The nonzero residual pairing}
Partition $G$ into target indices $A=(\lambda,\mu)$ and lower indices
$B=(\nu,\omega)$. The matrix $G_{BB}$ is positive definite, and
\[
 G_{AB}G_{BB}^{-1}
 =-\frac1{1513426}
     \begin{pmatrix}953923&134828\\292745&41134\end{pmatrix}.
\]
Thus the required polynomials are
\begin{align*}
 p_\lambda&=f_\lambda+
          \frac{953923 f_\nu+134828 f_\omega}{1513426},\\
 p_\mu&=f_\mu+
          \frac{292745 f_\nu+41134 f_\omega}{1513426}.
\end{align*}
By Sections 2--3 these lie in the claimed Jacobi spans and are
orthogonal to their full strictly lower spaces. Their normalized
Gram matrix is the Schur complement
\[
 a^{-6}\begin{pmatrix}
 \ip{p_\lambda}{p_\lambda}&\ip{p_\lambda}{p_\mu}\\
 \ip{p_\mu}{p_\lambda}&\ip{p_\mu}{p_\mu}
 \end{pmatrix}
 =G_{AA}-G_{AB}G_{BB}^{-1}G_{BA}
 =\frac1{24214816}
       \begin{pmatrix}3854156&-790\\-790&7201587\end{pmatrix}.
\]
The off-diagonal entry proves the theorem. If every pair of the
source's basis polynomials from the two distinct blocks were
orthogonal, every linear combination from the first would be
orthogonal to every linear combination from the second. The displayed
strictly negative number rules this out. The argument therefore does
not exploit the source's exception for equivalent indices.\qed

\section*{Task 2. Locate the smallest examples}
Determine the least number of matrix variables allowing failure at
$m=3$ and total degree six. In particular, investigate two variables
with $\tau_1=\tau_2=(2,1)$, whose product contains both target partition
types. For $m\le2$, the source's partition order is total, so distinct
non-equivalent indices are already comparable; this identifies a
natural boundary for the search.

\section*{Task 3. Classify the parameter dependence}
For six identically distributed matrix beta variables, replace the
ratio $b/a=-1/6$ by its general value. Determine the exact rational
function obtained from \eqref{gram} and the same Schur complement,
its zeros, and its sign throughout the admissible parameter region.
This should quantify how the obstruction persists under perturbation.

\section*{Task 4. Give individually labelled witnesses}
Fix explicit standard-tableau matrix units in each multiplicity space
and compute a particular pair of original Jacobi basis polynomials
with nonzero pairing. Compare its coefficients and size with the
basis-independent central certificate given here.

\section*{Task 5. Formulate a corrected orthogonality statement}
Determine which incomparable partition blocks must be grouped
together to admit an orthogonal construction preserving the original
individual matrix degrees. Study whether additional symmetric
operators can separate the resulting blocks, and compare this with
Gram--Schmidt for a total extension of dominance order.

\begin{thebibliography}{1}
\bibitem{K} T.~Kuijper,
\emph{Jacobi polynomials on the simplex of Hermitian matrices and
zonal spherical functions on partial flag manifolds},
\href{https://arxiv.org/abs/2607.25760v1}{arXiv:2607.25760v1} (2026),
equations (7), (8), (12), Definition 2.23 and Conjecture 2.25.
\end{thebibliography}
\noindent\textit{Verification and status.}
The accompanying Python verifier reproduces the finite arithmetic
certificate (SymPy and NumPy required).
The full current 40-page source retains the conjecture; dated claim
checks accompany this note. Mathematical writing and review were
performed by the authoring agent, without independent peer review.
\end{document}
